\color{unidarkgreen}{\subsection{Логика деления (ЛД)}}
\color{black}

\begin{enumerate}[label=\arabic{section}.\arabic{subsection}.\arabic*,  align=left, labelsep=4pt, leftmargin=0pt, itemindent=57pt]

\item
В составе логики деления предусмотрено действие от пуска ступеней ДЗ при междуфазных КЗ.

\begin{figure}[H]
\centering
\includegraphics[width=0.8\textwidth,height=0.8\textheight,keepaspectratio]{\fbpath/051.1791 ЛД/img/LD_DZ_FF_1.pdf}
\captionof{figure}{Функциональная схема ЛД от ДЗ-ФФ-1}\label{pic_ld:LD_DZ_FF_1}
\end{figure}

\item
Деление от ДЗ-ФФ-1(2…6) вводится в работу с помощью программного переключателя \textbf{«ЛД от ДЗ-ФФ-1(2…6)»}.

\item
При направлении ступени в сторону прилегающей сети, предусматривается возможность деления своей
стороны двумя действиями:
\begin{itemize}
\item Первым действием на предварительного деления отключение СВ с выдержкой времени \textbf{«Тср ДЗ-ФФ-1(2…6) на СВ»} или отключение выключателей выбранного сечения с выдержкой времени \textbf{«Тср ДЗ-ФФ-1(2…6) на В сеч»};
\item Вторым действием на отключение выключателей своей стороны с выдержкой времени \textbf{«Тср ДЗ-ФФ-1(2…6) на В»}.
\end{itemize}

\item
При направлении ступени в сторону автотрансформатора, предусматривается возможность пуска ЛД
смежной стороны с выдержкой времени \textbf{«Тср ДЗ-ФФ-1(2…6) на пуск ЛД см ст»} в резервные защиты смежной
стороны. В этом случае последовательность действия деления смежной стороны определяется защитами
смежной стороны.

\item
Ввод деления от ТНЗНП-1(2…6) и настройка параметров срабатывания осуществляется аналогично.

\item
В составе логики деления предусмотрено действие от пуска ступеней МТЗ.

\begin{figure}[H]
\centering
\includegraphics[width=0.8\textwidth,height=0.8\textheight,keepaspectratio]{\fbpath/051.1791 ЛД/img/LD_MTZ_1.pdf}
\captionof{figure}{Функциональная схема ЛД от МТЗ-1}\label{pic_ld:LD_MTZ_1}
\end{figure}

\item
Деление от МТЗ-1(2) вводится в работу с помощью программного переключателя \textbf{«ЛД от МТЗ-1(2)»}.

\item
При направлении ступени в сторону прилегающей сети, предусматривается возможность деления своей
стороны двумя действиями:
\begin{itemize}
\item Первым действием на предварительного деления отключение СВ с выдержкой времени \textbf{«Тср МТЗ-1(2) на СВ»} или отключение выключателей выбранного сечения с выдержкой времени \textbf{«Тср МТЗ-1(2) на В сеч»};
\item Вторым действием на отключение выключателей своей стороны с выдержкой времени \textbf{«Тср МТЗ-1(2) на В»}.
\end{itemize} 

\item
Ввод деления от ДЗ-ФЗ-1(2), от ОУ2 ДЗ, от ОУ2 ТНЗНП и настройка параметров срабатывания осуществляется аналогично.

\item
В составе логики деления предусмотрено действие от пуска ступеней ДЗ смежной стороны.

\begin{figure}[H]
\centering
\includegraphics[width=0.8\textwidth,height=0.8\textheight,keepaspectratio]{\fbpath/051.1791 ЛД/img/LD_DZ_sm.pdf}
\captionof{figure}{Функциональная схема ЛД от ДЗ смежной стороны}\label{pic_ld:LD_DZ_sm}
\end{figure}

\item
Деление от ДЗ смежной стороны вводится в работу с помощью программного переключателя \textbf{«ЛД от ДЗ см ст»}.

\item
Предусматривается возможность приема внешнего сигнала \textit{«П. ЛД от ДЗ см. ст»} от резервных защит
смежной стороны на формирование ЛД своей стороны двумя действиями:
\begin{itemize}
\item Первым действием на предварительного деления отключение СВ с выдержкой времени \textbf{«Тср ДЗ см ст на СВ»} или отключение выключателей выбранного сечения с выдержкой времени \textbf{«Тср ДЗ см ст на В сеч»}; 
\item Вторым действием на отключение выключателей своей стороны с выдержкой времени \textbf{«Тср ДЗ см ст на В»}.
\end{itemize}

\item
Действие ЛД от ДЗ смежной стороны контролируется пуском одной из ступеней ДЗ в сеть, выбор ступени
задается программным переключателем \textbf{«Контр ЛД от ДЗ см ст»}.

\item
Ввод ЛД от ТНЗНП смежной стороны и настройка параметров срабатывания осуществляется аналогично.

\begin{figure}[H]
\centering
\includegraphics[width=0.8\textwidth,height=0.8\textheight,keepaspectratio]{\fbpath/051.1791 ЛД/img/LD_GL.pdf}
\captionof{figure}{Функциональная схема выбора сечения}\label{pic_ld:LD_GL}
\end{figure}

\item
С помощью команды управления \textit{«Выбор сечения»} производится оперативный выбор предварительного
деления выключателей по 1 или 2 сечению. 

\item
Оперативный ввод/вывод логики деления производится с помощью БУ (блоков управления функциями) \textit{«Вывод ЛД»}. 

\small
\begin{longtable}{|>{\centering\arraybackslash}m{5.3cm}|>{\centering\arraybackslash}m{4.2cm}|>{\centering\arraybackslash}m{1cm}|>{\centering\arraybackslash}m{3.3cm}|>{\centering\arraybackslash}m{1.8cm}|}
% Первая страница: полный заголовок
\caption{Перечень уставочных параметров ЛД\hfill\vspace{-0.5\baselineskip}}\label{ld:tbl1}\\
\hline
\rowcolor{gray!30}
Наименование уставки & Диапазон & Шаг & Значение по умолчанию & Единица измерения \\
\hline
\endfirsthead

% Промежуточные страницы: «Продолжение таблицы»
\caption*{\hspace{3pt}\emph{Продолжение таблицы \ref{ld:tbl1}}\hfill\ }\vspace{-0.5\baselineskip} \\
\hline
\rowcolor{gray!30}
Наименование уставки & Диапазон & Шаг & Значение по умолчанию & Единица измерения \\
\hline
\endhead

%===m>ЛД|-
%===m
\end{longtable}
\normalsize	

\end{enumerate}